% Methods, left half, bottom: one real pixel as a tide gauge (step 4 of the
% flow chart with real data). Figure slot \figPixW x \figPixH (the delivered
% size: placed at 1:1); script tools/cg/fig_pixel_gauge.py (default preset
% "half"). What it is goes on the title line (\cgtitlenote: no height spent on
% a caption); the Copernicus credit is the footer's.
\setlength{\parskip}{0pt}%
\cgtitlenote{\cguse{\cgPixTitle}{A pixel is a tide gauge}}{\cguse{\cgPixNote}{one
  real pixel, R\'{\i}a de Villaviciosa, 2023--2025}}
\vspace{1.5mm}
\cgfig{pixel_gauge}{\figPixW}{\figPixH}{one real Villaviciosa pixel: every
  clear look, wet or dry, at its water level; the fitted flooding curve, its
  midpoint (the elevation) and its width (the uncertainty)}
